\sectionTitle{教育背景}{\faiconsixbf{graduation-cap}}
\datedsubsection{\texthl{哈尔滨工业大学} / \textit{应用统计} / \textit{硕士}}{2025 -- 2028}
\normalline{研究方向为统计建模与数据分析}
\vspace{-3pt}
\datedsubsection{\texthl{哈尔滨工业大学} / \textit{信息与计算科学} / \textit{本科}}{2020 -- 2024}
\normalline{数学与计算机交叉学科背景}

\sectionTitle{科研经历}{\faiconsixbf{flask}}
\datedsubsection{\textbf{体质数据分析研究} 哈尔滨}{2025.09 -- 至今}
\role{Python, PyTorch, 统计分析}{导师课题组}
对体质健康数据进行统计建模与分析
\begin{itemize}
  \item 主导数据处理、特征工程与统计建模全流程
  \item 应用多元统计方法探索体质影响因素
  \item 撰写学术论文并投稿《体育科学》
\end{itemize}

\sectionTitle{竞赛获奖}{\faiconsixbf{award}}
\begin{onehalfspacing}
\datedline{2026\space 全国统计建模大赛 / \textit{研究生组}}{2026}
\datedline{2023\space 全国大学生数学建模竞赛 / \textit{黑龙江省二等奖}}{2023}
\datedline{2022\space 美国大学生数学建模竞赛 / \textit{Meritorious Winner}}{2022}
\end{onehalfspacing}

\sectionTitle{专业技能}{\faiconsixbf{gears}}
\begin{onehalfspacing}
\normalline{熟练掌握 Python 数据分析与机器学习工具栈 (PyTorch, pandas, NumPy, scikit-learn)}
\normalline{熟悉 Linux/macOS 开发环境，熟练使用 Git 版本控制}
\normalline{具备 \LaTeX 学术排版能力，英语流利 (CET-6)}
\end{onehalfspacing}

\sectionTitle{其他经历}{\faiconsixbf{school}}
\begin{onehalfspacing}
\datedline{高中阶段参加全国信息学奥林匹克竞赛 (NOI)}{}
\datedline{入党积极分子}{}
\end{onehalfspacing}

\sectionTitle{自我评价}{\faiconsixbf{comment}}
\begin{onehalfspacing}
\normalline{数学基础扎实，编程能力强，善于用数据驱动的方式解决实际问题。对统计建模与机器学习有浓厚兴趣，具备独立科研与团队协作能力。}
\end{onehalfspacing}
